\documentclass[12pt]{article}
\usepackage[utf8]{inputenc}
\usepackage[T1]{fontenc}
\usepackage{hyperref}
\title{Ocotl\'an Zapotec Language Resources for K-12 Curriculum Development}
\author{Compiled by Claude Sonnet 5.5 (claude-sonnet-5-5) for Andr\'as Kornai}
\date{}
\begin{document}
\maketitle

\section{Introduction}
Ocotl\'an Zapotec (ISO 639-3 zac; Glottocode ocot1244) is an Otomanguean language of Oaxaca, Mexico, with about 15,000 speakers. Global Recordings Network lists several varieties, including Asunci\'on Ocotl\'an, Santa Luc\'ia and La Garzona. Searches found mainly religious and audio materials, and no school materials or dictionary specific to zac. Chichic\'apam Zapotec (zpv) is a different language and its resources are not cited here.

\section{Primary Educational Resources}
None identified in my searches.

\section{Linguistic and Reference Materials}
No grammar or dictionary specific to zac turned up in my searches. A multi-variety dictionary is listed under Digital and Audio Resources.

\section{Religious and Community Resources}
\begin{itemize}
\item Scripture Earth, Zapotec, Ocotl\'an (New Testament with audio and video): \url{https://www.scriptureearth.org/00eng.php?iso=zac}
\end{itemize}

\section{Digital and Audio Resources}
\begin{itemize}
\item Global Recordings Network, Zapotec, Ocotlan (audio-visual Bible lessons and Bible stories): \url{https://www.globalrecordings.net/en/language/zac}
\item Global Recordings Network, Asunci\'on variety: \url{https://www.globalrecordings.net/en/language/4950}
\item Dictionary of the Zapotec Language by \'Oscar M\'endez Espinoza (five volumes, many communities, not confirmed to include Ocotl\'an), hosted by AILLA: \url{https://islandora-ailla.lib.utexas.edu/islandora/object/ailla%3A279600}
\end{itemize}

\section{Curriculum Development Context}
The language has a substantial community, but documented teaching resources are scarce. The New Testament and audio recordings offer a base of written and spoken text in the language. A curriculum would need new primers, an agreed orthography and work with speakers of the different varieties.

\end{document}